%
% Front matter
%

\begin{frontmatter}

\title{
  T2K 1.431$\times10^{20}$-POT {\em 3-Flavour} Muon-Neutrino Disappearance Analysis
}

\vspace*{0.7cm}

\author[RAL]      {Costas Andreopoulos}
\author[RAL,OXF]  {Thomas Dealtry}
\author[VALENCIA] {Lorena Escudero S\'anchez}
\author[LANCS]    {and Nick Grant}

\address[RAL]{Particle Physics Dept., Rutherford Appleton Laboratory, STFC \\
              Harwell Science and Innovation Campus\\
              Oxfordshire OX11 0QX, UK}
\address[OXF]{Subdepartment of Particle Physics, University of Oxford \\
              Denys Wilkinson Building, Keble Road\\
              Oxfordshire OX1 3RH, UK}
\address[LANCS]{Physics Dept., Lancaster University \\
              Lancaster, LA1 4YB, UK}
\address[VALENCIA]{Instituto de Física Corpuscular, Edificio Institutos de Investigación \\
              Apartado de Correos 22085, E-46071 Valencia, Spain}

%
% Abstract
%

\begin{abstract}

We present the results of a 3-flavour $\nu_{\mu}$-disappearance analysis on the combined T2K Run 1+2 
dataset, which corresponds to an integrated J-PARC neutrino beam exposure of 1.431 $\times$ 10$^{20}$ POT. 
This analysis used 3-flavour oscillation probabilities instead of the 2-flavour approximation that was used 
in our earlier analysis reported in Ref.\ \cite{OurRun1and2Analysis2F}. Also oscillated $\nu_{e}$ templates 
were included in this analysis in the interaction modes CCQE, CC1$\pi$ and CC other, whereas these were 
neglected in the 2-flavour analysis. 

In the 3-flavour analysis, the total predicted number of single $\mu$-like ring events in Super-K is higher
for the inverted hierarchy (IH) than for the normal hierarchy (NH) if the atmospheric mass-squared difference
is $|\Delta m^{2}_{32}|$. This results from the prediction of more events above the oscillation maximum for 
the IH than for the NH and less below it, which is due to the 3-flavour $\nu_{\mu}$-survival probability being 
higher for the IH than the NH above the oscillation maximum and lower below the oscillation maximum. Since there 
are more events above the oscillation maximum than below it, the total predicted number is higher for the IH 
than for the NH. This is true for any combination of the non-23 oscillation parameters.

Uncertainties in $\delta_{CP}$ can change the predicted reconstructed spectrum of single $\mu$-like ring events
in Super-K by up to 5\% above the oscillation maximum and by up to 3\% below it. Meanwhile changes
in the spectrum due to uncertainties in $\sin^{2}(2 \theta_{13})$, $\sin^{2}(2 \theta_{12})$ and $\Delta m_{12}^{2}$ 
are relatively small at $\leq$ 0.5\%.

A 3-flavour $\nu_{\mu}$-disappearance fit was made of the observed Super-K reconstructed energy spectrum of
single $\mu$-like ring events in the combined Run 1+2 dataset. In this fit, sin$^{2}$(2$\theta_{23}$) and
$|\Delta m^{2}_{32}|$ were varied to find their best-fit values, but the non-23 oscillation parameters were
fixed. Using GENIE, this gave best-fit values of
sin$^{2}$(2$\theta_{23}$) = 0.98, $|\Delta m^{2}_{32}|$ = 2.6 $\times$ 10$^{-3}$ eV$^{2}$/c$^{4}$ for both mass
hierarchies, while a fit using NEUT gave sin$^{2}$(2$\theta_{23}$) = 0.99, $|\Delta m^{2}_{32}|$
= 2.6 $\times$ 10$^{-3}$ eV$^{2}$/c$^{4}$ for the NH and sin$^{2}$(2$\theta_{23}$) = 0.99, $|\Delta m^{2}_{32}|$
= 2.7 $\times$ 10$^{-3}$ eV$^{2}$/c$^{4}$ for the IH. For comparison, the best-fit values from the 2-flavour
analysis in \cite{OurRun1and2Analysis2F} were sin$^{2}$(2$\theta_{23}$) = 0.97, $|\Delta m^{2}_{32}|$
= 2.6 $\times$ 10$^{-3}$ eV$^{2}$/c$^{4}$ using GENIE and sin$^{2}$(2$\theta_{23}$) = 0.98, $|\Delta m^{2}_{32}|$
= 2.6 $\times$ 10$^{-3}$ eV$^{2}$/c$^{4}$ using NEUT. The 3-flavour best-fit points were higher in sin$^{2}$(2$\theta_{23}$)
than those in the 2-flavour analysis due to the effect of sin$^{2}$(2$\theta_{13}$): the 3-flavour analysis
used the value sin$^{2}$(2$\theta_{13}$) = 0.1, whereas the 2-flavour analysis used the approximation 
sin$^{2}$(2$\theta_{13}$) = 0.

Uncertainties in the best-fit points due to uncertainties in $\delta_{CP}$ were larger than those due to
uncertainties in $\sin^{2}(2 \theta_{13})$, $\sin^{2}(2 \theta_{12})$ and $\Delta m_{12}^{2}$; these uncertainties 
due to $\delta_{CP}$ were 0.5\% (NH) and 0.35\% (IH) in sin$^{2}$(2$\theta_{23}$) and 0.6\% (NH) and 0.4\% (IH) in
$|\Delta m^{2}_{32}|$. However this is relatively small compared with the combined systematic uncertainties in
the best-fit points which were 4.3 - 5.1\%.

The 3-flavour confidence contours are narrower in sin$^{2}$(2$\theta_{23}$) than the 2-flavour contours
since the 3-flavour best-fit points are nearer to sin$^{2}$(2$\theta_{23}$) = 1.0. The confidence contours
for the IH include higher values of $|\Delta m^{2}_{32}|$ than those for the NH; this is because it is necessary
to choose a higher value of this parameter for the IH than for the NH in order to fit the oscillation dip at a
certain energy in the real data.

\end{abstract}

%
% Keywords
%

%\begin{keyword}
%\end{keyword}

\end{frontmatter}

\clearpage

\tableofcontents

\clearpage

\listoffigures
\listoftables


